%!TEX root = thesis.tex

\chapter{Introduction}
\label{ch:introduction}

\section{Motivation}
\label{sec:intro:motivation}

\ac{CRC} is the third most commonly diagnosed cancer and the second leading cause of cancer-related mortality worldwide~\cite{sung2021global}. The detection and removal of precancerous polyps during colonoscopy has been shown to significantly reduce \ac{CRC} incidence and mortality~\cite{zauber2012colonoscopic}. However, the polyp miss rate during routine colonoscopy remains alarmingly high, with studies reporting that up to 26\% of polyps may be overlooked~\cite{kim2017miss}. This motivates the development of computer-aided detection and segmentation systems that can assist endoscopists in identifying and delineating polyps in real time.

Deep learning approaches based on \acp{CNN} have achieved considerable success in automated polyp segmentation~\cite{fan2020pranet, jha2020doubleunet, zhou2018unetplusplus}. However, these purely vision-based models treat segmentation as an isolated pixel classification task, ignoring the rich clinical context that endoscopists routinely leverage---such as polyp morphology (pedunculated vs.\ sessile), size, anatomical location, and surface characteristics. This gap between algorithmic processing and clinical reasoning limits both the accuracy and interpretability of existing systems.

\section{Problem Statement}
\label{sec:intro:problem}

This Praktikum investigates the hypothesis that \textbf{semantic language cues improve both segmentation quality and interpretability of uncertainty in medical image segmentation}. Specifically, we address three research questions:

\begin{enumerate}
	\item Does integrating structured clinical text descriptions with visual features improve polyp segmentation accuracy compared to a vision-only baseline?
	\item Can language-guided models produce better-calibrated uncertainty estimates, particularly in clinically relevant regions such as polyp boundaries?
	\item Which clinical attributes (shape, size, location, pathology) contribute most significantly to segmentation improvement?
\end{enumerate}

\section{Contributions}
\label{sec:intro:contributions}

The main contributions of this work are:

\begin{itemize}
	\item \textbf{A modular \ac{VL} segmentation framework} that combines a \ac{ResNet}-34 vision encoder with a frozen \ac{PubMedBERT} text encoder through cross-attention fusion, enabling integration of structured clinical descriptions into the segmentation pipeline.
	
	\item \textbf{A structured caption generation system} based on the Paris Classification that automatically generates clinically grounded textual descriptions from polyp metadata (morphology, size, location, and surface characteristics).
	
	\item \textbf{Three novel uncertainty evaluation metrics}:
	\begin{itemize}
		\item \ac{URS}: Quantifies the reduction in predictive uncertainty when language cues are incorporated.
		\item \ac{SAS}: Measures the correlation between text-derived semantic similarity and spatial uncertainty patterns.
		\item \ac{AC-ECE}: Evaluates calibration quality conditioned on specific clinical attributes.
	\end{itemize}
	
	\item \textbf{A comprehensive ablation study} across six experimental configurations that isolates the contribution of individual clinical attributes to segmentation quality and uncertainty estimation.
\end{itemize}

\section{Report Structure}
\label{sec:intro:structure}

The remainder of this report is organized as follows:

\begin{itemize}
	\item \textbf{Chapter~\ref{ch:background}} provides the necessary background on medical image segmentation, vision-language models, and uncertainty estimation.
	\item \textbf{Chapter~\ref{ch:relatedwork}} reviews related work in polyp segmentation, vision-language medical imaging, and uncertainty quantification.
	\item \textbf{Chapter~\ref{ch:approach}} details the proposed architecture including the vision encoder, text encoder, fusion mechanisms, decoder, and uncertainty estimation framework.
	\item \textbf{Chapter~\ref{ch:implementation}} describes the software implementation, data pipeline, and training procedure.
	\item \textbf{Chapter~\ref{ch:evaluation}} presents the experimental setup, evaluation metrics, and results.
	\item \textbf{Chapter~\ref{ch:conclusion}} summarizes the findings and discusses future work.
\end{itemize}
